\documentclass[11pt, a4paper]{article}

\usepackage[T1]{fontenc}
\usepackage[utf8]{inputenc}
\usepackage{tgpagella}
\usepackage[spanish, es-noshorthands]{babel}
\usepackage{amsmath, amssymb, bm}
\usepackage{tabularx}
\usepackage{booktabs}
\usepackage[most]{tcolorbox}
\usepackage[margin=2.5cm]{geometry}
\usepackage{ragged2e}
\usepackage{fancyhdr}
\usepackage{tikz}

\pagestyle{fancy}
\fancyhf{}
\setlength{\headheight}{16pt}
\lhead{\textbf{UCB Sede Tarija} | Ing. Mecatrónica}
\chead{}
\rhead{IMT-121 Circuitos Electrónicos I | Síntesis y Puente EC3}
\lfoot{Prof: M.Sc. Ing. Bernardo Quiroga Turdera}
\rfoot{Página \thepage}
\renewcommand{\headrulewidth}{0.4pt}

\definecolor{ucbblue}{RGB}{0, 51, 102}

\begin{document}

\begin{tcolorbox}[colback=ucbblue!5!white, colframe=ucbblue, title=\textbf{Síntesis de Circuitos DC (EC2) y el Puente hacia Transitorios (EC3)}]
\end{tcolorbox}

\section*{1. Lo que EC2 nos entregó (El Modelo Mental Estático)}
A lo largo de esta unidad hemos demostrado que cualquier red resistiva lineal, sin importar su complejidad interna, puede ser abstraída y validada mediante la siguiente cadena de razonamiento:
\begin{center}
    \textbf{Red $\rightarrow$ Superposición $\rightarrow$ Puerto $\rightarrow$ Equivalente (Thévenin/Norton) $\rightarrow$ Carga $\rightarrow$ Máx. Potencia $\rightarrow$ Validación}
\end{center}
En este modelo (DC Estático), las relaciones son puramente \textbf{algebraicas} ($V = I \cdot R$). Si cambiamos un interruptor, el voltaje y la corriente cambian instantáneamente en el mismo instante $t$.

\section*{2. El Límite de EC2: Lo que aún no podemos describir}
¿Qué ocurre físicamente en un circuito en el instante exacto en que accionamos un interruptor para encender una máquina? ¿El voltaje y la corriente llegan a su valor final en cero segundos?
\begin{itemize}
    \item Los circuitos resistivos \textbf{no almacenan energía}.
    \item Los \textbf{Capacitores} y los \textbf{Inductores} sí lo hacen. Esto introduce una variable que hasta ahora ignorábamos: \textbf{El Tiempo ($t$)}.
\end{itemize}

\section*{3. Principios de Continuidad (Preview EC3)}
En la próxima unidad estudiaremos cómo los circuitos evolucionan desde una condición inicial hacia un estado final (Transitorio). Existen dos reglas de oro inviolables en la física de circuitos:

\vspace{0.2cm}
\noindent\begin{minipage}{0.48\textwidth}
    \textbf{El Capacitor ($C$): Inercia de Voltaje} \\
    El voltaje en los terminales de un capacitor ideal \textbf{no puede cambiar instantáneamente}.
    \begin{equation*}
        v_C(0^+) = v_C(0^-)
    \end{equation*}
    La corriente depende de la \textit{derivada} del voltaje temporal: $i_C = C \frac{dv_C}{dt}$.
\end{minipage}
\hfill
\begin{minipage}{0.48\textwidth}
    \textbf{El Inductor ($L$): Inercia de Corriente} \\
    La corriente que atraviesa un inductor ideal \textbf{no puede cambiar instantáneamente}.
    \begin{equation*}
        i_L(0^+) = i_L(0^-)
    \end{equation*}
    El voltaje depende de la \textit{derivada} de la corriente temporal: $v_L = L \frac{di_L}{dt}$.
\end{minipage}

\vspace{0.5cm}
\begin{tcolorbox}[colback=gray!5, colframe=gray!50, title=\textbf{Reflexión Conceptual}]
\textbf{Pregunta Clave:} Si las relaciones constitutivas de $C$ y $L$ involucran derivadas respecto al tiempo ($\frac{d}{dt}$), ¿qué tipo de ecuaciones matemáticas creen que tendremos que resolver para encontrar $v(t)$ e $i(t)$ a partir de la próxima semana, abandonando el álgebra simple de las resistencias?
\end{tcolorbox}

\end{document}
